\section{A full positive energy calculation in a cylindrical benchmark}

The following calculation signs the entire ideal-MHD quadratic energy,
including the interface curvature term.  It is a benchmark on a periodic
Euclidean cylinder, not a claim that a periodic cylinder is an embedded
stellarator in three-dimensional Euclidean space.  Its useful feature is
that the estimates are uniform in the axial period.

Let the plasma occupy the solid periodic cylinder (including its axis,
where all fields and displacements are required to be smooth)
\[
 r<a,\qquad \theta\in\mathbb R/(2\pi\mathbb Z),
 \qquad z\in\mathbb R/(L\mathbb Z),
\]
and let the vacuum occupy $a<r<b$.  The outer surface $r=b$ is fixed and
perfectly conducting.  Take the equilibrium fields
\[
 Y=B_i e_z,\qquad
 X=B_e e_z+B_\theta\frac{a}{r}e_\theta,
 \qquad B_i>0.
\]
Both fields are curl-free and divergence-free in their respective open
regions.  The plasma pressure and density are positive constants.  The
interface has the surface current required by the jump of the tangential
magnetic field.  Force balance is
\[
 p=\frac{B_e^2+B_\theta^2-B_i^2}{2\mu_0}.
\]

We use the lossless ideal-MHD energy principle.  In particular, the outer
wall is passive, there is no work from a perturbed external power supply,
and the magnetic fluxes through the independent material cuts are fixed.
One convenient complete choice fixes the plasma axial flux, the vacuum
axial flux, and the vacuum azimuthal flux.  The last two conditions specify
the harmonic parts of the vacuum perturbation.  The condition that the
outer wall is fixed is not, by itself, a specification of these harmonic
parts.  The calculations below also remain lower bounds if additional
admissible harmonic vacuum perturbations are retained: their magnetic
energy is nonnegative and orthogonal to the gradient part.

Write $\eta=\xi\cdot e_r$ for the interface displacement and
\[
 Q=\nabla\times(\xi\times Y)
\]
for the plasma magnetic perturbation.  If $b_v$ is the vacuum magnetic
perturbation, the complete quadratic energy is
\[
\begin{split}
2\delta W={}&
 \int_{r<a}\left(\frac{|Q|^2}{\mu_0}
       +\gamma p\,|\operatorname{div}\xi|^2\right)\,dx
 +\frac{1}{\mu_0}\int_{a<r<b}|b_v|^2\,dx
 -\frac{B_\theta^2}{\mu_0 a}\int_{r=a}|\eta|^2\,dS.
\end{split}
\]
Here $\gamma>1$ is the adiabatic exponent.  In particular, the negative
surface term has not been omitted.  Its coefficient follows from
\[
 C=\frac{\mathrm{II}(Y,Y)-\mathrm{II}(X,X)}{\mu_0}
   =-\frac{B_\theta^2}{\mu_0a},
\]
where $\mathrm{II}(v,w)=v\cdot\nabla_w n$ and $n=e_r$.

\subsection{Exact mode minimization}

Expand
\[
 \eta(\theta,z)=\sum_{m,n\in\mathbb Z}
       \eta_{mn}\exp(i m\theta+i k_nz),
 \qquad k_n=\frac{2\pi n}{L}.
\]
The normalization is
\[
 \frac{1}{2\pi aL}\int_{r=a}|\eta|^2\,dS
   =\sum_{m,n}|\eta_{mn}|^2.
\]
Put $\kappa=a|k|$.  Derivatives of $I_m$ and $K_m$ below mean derivatives
with respect to their arguments.

First suppose $k\ne0$, and replace $m$ by $|m|$ when defining the radial
functions.  If $w$ is the transverse part of $\xi$, then
\[
 Q_\perp=i kB_iw,\qquad Q_z=-B_i\operatorname{div}_\perp w.
\]
Consequently
\[
 \int |Q|^2\,dx
  =B_i^2\int\left(k^2|w|^2+
                         |\operatorname{div}_\perp w|^2\right)\,dx.
\]
The minimum for prescribed $w\cdot e_r=\eta$ is attained by
$w=\nabla_\perp\chi$, where
\[
 (\Delta_\perp-k^2)\chi=0,
 \qquad \partial_r\chi(a)=\eta.
\]
Indeed, writing $w=\nabla_\perp\chi+v$ with $v\cdot e_r=0$, the cross term
vanishes by integration by parts.  The minimizing longitudinal
displacement is $\xi_z=i k\chi$, which gives
$\operatorname{div}\xi=0$.  Thus dropping the positive compressibility
term loses nothing for these nonzero-$k$ modes.  The plasma contribution
per unit interface area and per $|\eta|^2$ is
\[
 \frac{B_i^2}{\mu_0a}
       \frac{\kappa I_m(\kappa)}{I_m'(\kappa)}.
\]

At the inner boundary the vacuum condition is
\[
 b_v\cdot e_r=\operatorname{div}_{\Sigma}(\eta X)
       =i\left(kB_e+\frac{mB_\theta}{a}\right)\eta,
\]
and $b_v\cdot e_r=0$ at $r=b$.  For a nonconstant Fourier mode, the vacuum
perturbation is the gradient of a single-valued harmonic potential.  Let
$N_{\rm out}(m,k)>0$ be its Neumann impedance, so that its energy per unit
inner-boundary area is
\[
 N_{\rm out}(m,k)
       \left|kB_e+\frac{mB_\theta}{a}\right|^2|\eta|^2.
\]
Set $n_{\rm out}=N_{\rm out}/a$.  The exact minimized coefficient is
\[
 \frac{2\delta W_{mn}}
      {(2\pi aL)|\eta_{mn}|^2}
 =\frac{B_i^2}{\mu_0a}\,F_{mn},
\]
where, for $k\ne0$,
\[
 F_{mn}=
 \frac{\kappa I_m(\kappa)}{I_m'(\kappa)}
 +\left(\frac{k aB_e+mB_\theta}{B_i}\right)^2 n_{\rm out}
 -\frac{B_\theta^2}{B_i^2}.
\]
For arbitrary plasma displacements the right side is a lower bound;
the omitted difference is nonnegative.

For $k=0$ and $m\ne0$, a transverse divergence-free displacement gives
zero plasma magnetic and compression energy.  Direct solution of the
vacuum problem gives
\[
 N_{\rm out}(m,0)=\frac{a}{|m|}
        \coth\left(|m|\log\frac ba\right),
\]
and hence
\[
 F_{m0}=\frac{B_\theta^2}{B_i^2}
     \left(|m|\coth\left(|m|\log\frac ba\right)-1\right)>0.
\]
In an unbounded exterior this becomes
$(|m|-1)B_\theta^2/B_i^2$: the negative curvature term cancels exactly
for a rigid transverse translation, $|m|=1$.  Thus a negative curvature
coefficient alone is not an instability proof.  The fixed close wall
removes this particular zero mode.

\subsection{An exact unstable branch, even with a close wall}

The choice of pressure-lifting direction matters.  Suppose $B_e\ge B_i$
and $B_\theta\ne0$.  For $m=1$, choose
\[
 k=-\frac{B_\theta}{aB_e}.
\]
The vacuum normal perturbation then vanishes identically, independently
of the wall distance.  With $\kappa=|B_\theta|/B_e$, the identity
$\kappa I_1'=I_1+\kappa I_2$ gives
\[
 \frac{\kappa I_1}{I_1'}<\kappa^2.
\]
Consequently the exact energy coefficient is strictly less than
\[
 \frac{B_\theta^2}{B_e^2}
       -\frac{B_\theta^2}{B_i^2}\le0.
\]
This is a genuine negative-energy mode whenever the indicated $k$ is
allowed by the axial period.  By continuity, every sufficiently long
period admits a nearby unstable mode.  Thus positive pressure,
noncollinearity, and a close wall do not stabilize arbitrary choices
of the two tangential fields.  The stable family below uses a slightly
smaller exterior axial field, while keeping the total exterior magnetic
pressure larger than the interior value.

\subsection{A low-pressure, all-mode stability window}

Here is one explicit sufficient window.  Choose
\[
 0<s\le\frac14,\qquad
 B_\theta^2=sB_i^2,\qquad
 B_e^2=\left(1-\frac{3s}{4}\right)B_i^2.
\]
Both axial fields are nonzero, and
\[
 p=\frac{sB_i^2}{8\mu_0}>0,
 \qquad \frac{2\mu_0p}{B_i^2}=\frac{s}{4}.
\]
Let
\[
 h=\frac ba,\qquad \ell=\log h,
 \qquad D=\frac{h^2-1}{2},
 \qquad \beta_e=|B_e|/B_i=\sqrt{1-\frac{3s}{4}}.
\]
Assume
\[
 \ell+\frac{sD}{\beta_e^2}\le\frac{s}{8}.
\]
A simple sufficient choice is
\[
 1<h\le1+\frac{s}{16}.
\]
This is a close ideal conducting wall, not a distant coil condition.

For $m\ge1$, the positive series for $I_m$ gives
\[
 \frac{I_{m+1}(\kappa)}{I_m(\kappa)}
       \le\frac{\kappa}{2(m+1)}.
\]
Using $\kappa I_m'=mI_m+\kappa I_{m+1}$, we obtain
\[
 \frac{\kappa I_m}{I_m'}
 \ge \frac{\kappa^2}
             {m+\kappa^2/(2(m+1))}.
\]
There is also an elementary lower bound for the vacuum impedance:
\[
 n_{\rm out}\ge\frac{1}{m^2\ell+\kappa^2D}.
\]
To see this, use a radially constant potential as a trial function in
the dual variational characterization of the Neumann energy.  Its
energy per unit inner-boundary area, divided by the squared boundary
value, is $(m^2\ell+\kappa^2D)/a$.  Optimizing its amplitude proves the
displayed lower bound.

The least favorable helicity has vacuum factor
$(m\sqrt{s}-\kappa\beta_e)^2$.  Put
\[
 t=\frac{\kappa\beta_e}{m\sqrt{s}}\ge0.
\]
The preceding estimates give
\[
 \frac{F_{mn}}{s}\ge
 \frac{t^2}{\beta_e^2+st^2/2}
 +\frac{(1-t)^2}{\ell+st^2D/\beta_e^2}-1.
\]
For $0\le t\le1$, the elementary Cauchy inequality
\[
 \frac{t^2}{A}+\frac{(1-t)^2}{B}\ge\frac{1}{A+B}
 \qquad (A,B>0)
\]
and the wall assumption imply
\[
 \frac{F_{mn}}{s}\ge
 \frac{1}{1-s/8}-1\ge\frac{s}{8}.
\]
For $t\ge1$, the first term alone yields
\[
 \frac{F_{mn}}{s}\ge
 \frac{1}{\beta_e^2+s/2}-1
 =\frac{1}{1-s/4}-1\ge\frac{s}{4}.
\]
Therefore
\[
 F_{mn}\ge\frac{s^2}{8}
 \qquad\text{for every }m\ne0\text{ and every }k.
\]
The $k=0$ case can either be read directly from the exact formula or
obtained as the limit of this estimate.

For $m=0$, $k\ne0$, the series inequality $I_1\le\kappa I_0/2$ gives
\[
 F_{0n}\ge\frac{\kappa I_0}{I_1}-s\ge2-s>0.
\]
The vacuum term is nonnegative and has not been needed for this bound.

For completeness, the stated simple wall choice implies the wall
assumption.  Indeed, $\ell\le s/16$, $\beta_e^2\ge13/16$, and
\[
 D\le\frac{s}{16}+\frac{s^2}{512}
     \le\frac{129s}{2048}.
\]
It follows that
\[
 \ell+\frac{sD}{\beta_e^2}
 \le\left(\frac1{16}+\frac{129}{6656}\right)s
 <\frac{s}{8}.
\]

\subsection{The spatially constant displacement and the flux constraints}

For $(m,k)=(0,0)$, it is important not to discard the plasma flux
constraint.  A uniform displacement $\eta$ requires
\[
 \delta B_i=-\frac{2B_i\eta}{a},
 \qquad \operatorname{div}\xi=\frac{2\eta}{a}
\]
for the energy-minimizing radial displacement.  Thus the plasma
magnetic and gas contributions per unit interface area are
\[
 \frac{2B_i^2}{\mu_0a}|\eta|^2
       +\frac{2\gamma p}{a}|\eta|^2.
\]
Even if the nonnegative vacuum contribution is discarded, the full
coefficient satisfies
\[
 F_{00}\ge2-s+\frac{\gamma s}{4}>0.
\]

The harmonic vacuum terms can also be written explicitly.  Fixing the
vacuum axial flux and azimuthal flux gives
\[
 \delta B_e=\frac{B_e\eta}{aD},
 \qquad
 \delta(aB_\theta)=\frac{B_\theta\eta}{\ell}.
\]
Their energies add respectively $\beta_e^2/D$ and $s/\ell$ to
$F_{00}$.  Thus in this flux ensemble
\[
 F_{00}=2-s+\frac{\gamma s}{4}
                +\frac{\beta_e^2}{D}+\frac{s}{\ell}.
\]
These contributions are orthogonal to the nonconstant potential
modes.  A calculation that specified only the outer wall and silently
set all magnetic periods to zero would miss these constraints.

\subsection{An explicit uniform interface gap}

The preceding positive constant bound can be strengthened to a
half-derivative interface estimate.  For $m\ge1$, the integral
representations of the Bessel functions imply
\[
 \frac{I_{m+1}(\kappa)}{I_m(\kappa)}<1,
 \qquad \frac{K_{m-1}(\kappa)}{K_m(\kappa)}\le1.
\]
The first ratio is the expectation of a number in $(-1,1)$ in the
positive integral representation of $I_m$; the second follows from
$\cosh((m-1)t)\le\cosh(mt)$ in the integral representation of $K_m$.
They give
\[
 \frac{\kappa I_m}{I_m'}\ge\frac{\kappa^2}{m+\kappa},
 \qquad n_{\rm out}\ge\frac{1}{m+\kappa}.
\]
For the second inequality, the wall impedance is at least the
unbounded-exterior impedance.  This follows by restricting trial
functions in the dual Neumann variational principle to the smaller
annulus.  The unbounded-exterior value is
$-K_m/(\kappa K_m')$, and the displayed ratio bound applies.

The quadratic form
\[
 \kappa^2+(m\sqrt{s}-\kappa\beta_e)^2
\]
has determinant $s$ and trace
$s+1+\beta_e^2=2+s/4\le33/16$.  Its least eigenvalue is therefore at
least $16s/33$.  Consequently, with $x=m+\kappa$,
\[
 F_{mn}\ge\frac{8s}{33}x-s.
\]
If $x\ge9$, this is at least $4sx/33$.  If $x\le9$, the previously
proved estimate $F_{mn}\ge s^2/8$ applies.  Combining the two gives
the convenient conservative bound
\[
 F_{mn}\ge\frac{s^2}{80}(1+m+\kappa)
 \qquad (m\ge1).
\]
For $m=0$, use $I_1<I_0$ in addition to the earlier series bound to
obtain $F_{0n}\ge\max(2-s,\kappa-s)$; this is more than sufficient for
the same estimate.  The zero mode was checked separately above.

We have therefore proved, for every smooth admissible displacement,
\[
 2\delta W\ge
 \frac{2\pi aL\,B_i^2}{\mu_0a}\,
 \frac{s^2}{80}
 \sum_{m,n}\left(1+|m|+a|k_n|\right)|\eta_{mn}|^2.
\]
The constant is independent of $L$.  This is a genuine positive
interface $H^{1/2}$ energy gap for the complete quadratic functional.
The nonnegative bulk excess over the mode minima is still present.
There remain bulk relabelling directions with $\eta=0$, $Q=0$ and
$\operatorname{div}\xi=0$.  The theorem does not give a coercive norm
of every material displacement, a nonlinear continuation theorem, or
a kinetic confinement theorem.

\section{A geometric sufficient condition on an actual toroidal interface}

Use the same passive, lossless exterior and fixed-flux ensemble as above.
Let $Y$ be a tangent vacuum field on an analytic toroidal interface,
straightened to $Y=\omega^i\partial_{\theta_i}$, and let $g$ and
$\mathrm{II}$ be the induced first and second fundamental forms.
For the pressure-lifting Hamiltonian
\[
 H(\theta,I)=\omega\cdot I+\frac12 I^{\mathsf T}g^{-1}(\theta)I,
\]
a selected family of KAM graphs with
\[
 I=r d+o(r)
\]
gives exterior trace $X=Y+r g^{-1}d+o(r)$.  A sufficient condition for
positive pressure and a positive full surface energy term is
\[
 \omega\cdot d>0,
 \qquad \mathrm{II}(Y,g^{-1}d)<0\quad\hbox{everywhere on the interface}.
\]
Strictness and compactness give a uniform margin.  For sufficiently
small selected $r>0$,
\[
 p>0,\qquad
 C=\frac{\mathrm{II}(Y,Y)-\mathrm{II}(X,X)}{\mu_0}>0.
\]
Every term in the complete ideal-MHD quadratic energy is then
nonnegative, including the interface term.  This is a sufficient
linear energy-stability criterion; it does not assume the positivity
that was missing from the earlier calculation.

Its feasibility is explicitly one-dimensional.  Normalize
$\omega\cdot d=1$ and write
\[
 d=\frac{\omega}{|\omega|^2}+\sigma J\omega,
\]
where $J$ is rotation by $\pi/2$ in the two coordinate components.
Define
\[
 a_0(\theta)=\frac{\mathrm{II}(Y,g^{-1}\omega)}{|\omega|^2},
 \qquad b_0(\theta)=\mathrm{II}(Y,g^{-1}J\omega).
\]
The complete condition is simply
\[
 a_0(\theta)+\sigma b_0(\theta)<0\quad\hbox{for all }\theta.
\]
Equivalently, $\sigma$ must lie strictly above
$\sup_{b_0<0}(-a_0/b_0)$ and strictly below
$\inf_{b_0>0}(-a_0/b_0)$; every point with $b_0=0$ must have $a_0<0$.
Empty sets impose no bound.  Compactness supplies the uniform strict
margin once such a $\sigma$ exists.  Failure has a one-point or a
two-point witness.  This is a concrete geometric design test.

If noncollinearity is separately desired, one adds either sign of
\[
 \det(\omega,g^{-1}d)\ne0
\]
uniformly on the torus.  These are additional affine inequalities in
$\sigma$.  Positivity of $C$ itself does not imply noncollinearity;
noncollinearity is not needed merely to sign the displayed quadratic
energy when $C>0$ has already been proved.

\section{What remains before transferring the positive benchmark to a torus}

The cylindrical theorem removes a real difficulty: it gives an
interface gap uniform in the axial period, rather than a calculation
restricted to finitely many modes.  It does not, by itself, construct
the requested toroidal equilibrium.

There is an exact variational reduction that could support such a
transfer.  On a specified solid torus, the harmonic Neumann magnetic
field of fixed flux minimizes magnetic energy among all divergence-free
tangent fields with that flux.  If $B_h$ is this harmonic field, Hodge
orthogonality gives
\[
 \int |B|^2=\int|B_h|^2+\int|B-B_h|^2.
\]
At fixed total mass and fixed uniform specific entropy, Jensen's
inequality similarly minimizes the gas internal energy at uniform
density.  Thus the sum of the harmonic magnetic energies, at the
specified plasma and vacuum fluxes, and the uniform adiabatic gas
energy is an exact lower bound for the ideal potential energy on a
given plasma domain.  Its first shape derivative is the total-pressure
jump.  At the cylinder, the preceding mode calculation is its
interface Hessian, including the flux-dependent zero mode.

This suggests solving the Bernoulli shape equation by an elliptic
implicit-function theorem and preserving the positive Hessian.  A
bent tube of major radius $1/\varepsilon$ has, in cross-section and
arclength coordinates, the Euclidean metric
\[
 g_\varepsilon=dx^2+dy^2+(1+\varepsilon x)^2ds^2,
 \qquad L=\frac{2\pi}{\varepsilon}.
\]
For each fixed $L$, an ordinary local perturbation argument is not
enough: its permitted metric perturbation could depend on $L$.
The actual Euclidean closing condition ties $\varepsilon$ to $L$.
One must prove uniform estimates for the nonlinear harmonic-field
shape map and its inverse, including the magnetic-period constraints,
as $L$ increases.  Low axial Neumann modes make a naive estimate for
the scalar potentials singular.  Estimates for the fields and the
flux-constrained shape operator may cancel that singularity, but that
cancellation must be proved, not presumed from the interface spectrum.

A precise sufficient transfer certificate is a quadratic-form bound:
after identifying the boundaries and their flux constraints, the
change in the reduced interface Hessian must have operator norm less
than half the cylinder gap from $H^{1/2}$ to $H^{-1/2}$.  This would
preserve positivity of the whole energy because the bulk terms remain
nonnegative for a current-free core with constant pressure.  Such a
uniform bound and an exact toroidal Bernoulli solution have not been
established in this calculation.

There are further independent requirements.  A nonaxisymmetric
perturbation need not give a nonzero rotational transform, nested
magnetic surfaces, or exact quasisymmetry automatically.  The close
perfectly conducting wall and the edge current in the benchmark are
substantial assumptions.  Finite wall resistivity and kinetic wall
losses are absent.  Thus the positive result is a complete ideal-MHD
linear energy theorem for the stated benchmark and a concrete route
to a geometric extension, not a proof of hot stable fusion plasma.
